\section{Situated Urban and Societal Substrate}
\label{sec:societal-substrate}

GAWorld places persistent agents in a shared substrate whose components read
and write overlapping state.  The city constrains movement, local occupancy
changes perception, relationships affect encounters, and financial or life
events alter later choices.  This coupling supports questions that cross
individual and aggregate levels, but it also means that mechanism provenance
and configuration are part of every experimental treatment.

\subsection{City graph and local physical conditions}

The city map parses places, roads, transit lines, interiors, and land-use
attributes into a graph.  Nodes carry categories, coordinates, capacity,
opening hours, popularity, density, and display metadata.  Destination
resolution can match activity terms to place categories, while graph routing
and transport-mode rules estimate distance, time, and fare.  Mode speeds,
fixed delays, road-class multipliers, rush-hour penalties, area price levels,
and weather preferences are explicit parameters.  They make travel costs
inspectable; their presence does not imply transport-model calibration.

At each tick, the local-physical layer recomputes occupancy from stationary
agents, clears stale counts, and exposes the current location's crowding,
opening status, and weather.  A crowd anomaly is flagged when occupancy is both
high and sharply above the previous tick.  These snapshots can enter
perception and interruption generation, linking collective movement to
individual behavior.  The mechanism is deterministic given agent locations
and parameters, although the locations themselves may follow stochastic or
LLM-mediated decisions.

The broader environment system supports configured natural, economic,
political, technological, and intraday events.  In rule mode it samples from
configured event distributions using a local random generator.  In LLM mode a
model proposes a daily summary, events, and intraday rules that are parsed into
the same bounded event schema; configured fallback generation remains
available if that call fails.  Severity and anomaly tagging are rule-based.
Reports must therefore identify the generator mode rather than treating all
environmental histories as equivalent.

\subsection{Social ties and life events}

Relationships are typed records rather than undifferentiated neighbor edges.
Role templates specify channels, obligation, protection, and daily decay;
interaction signals update closeness and contact time.  The implementation can
bootstrap off-screen contacts from an LLM-generated backstory, with a
deterministic heuristic fallback.  It can also generate events involving
off-screen contacts, using templates and optional LLM wording, and disclose
shared contacts across agents.

Maintenance is rule-based.  Uncontacted ties decay at role-specific rates,
long-neglected obligations can rise, and an enforcement pass assigns nested
Dunbar-style tiers while pruning beyond a configured outer cap and protecting
selected kin roles.  These tiers are an operational capacity constraint, not
evidence that the simulated network reproduces measured human networks.

Life events are normalized, persisted, scheduled by simulated day and time,
and drained once when due.  Events can target agents and carry bounded state
effects and tags.  This provides a controlled exogenous-event surface whose
timing and direct updates are inspectable.  The realism of event frequencies,
dependencies, and consequences must be established by a separate study.

\subsection{Economic state transitions}

The finance module maintains checking, savings, investment, and housing-fund
accounts; monthly budgets; income and expense records; liabilities and credit;
sector pools; government flows; and macro state.  Configured rules calculate
social-insurance deductions and progressive income tax, route wages and
consumption, enforce available-cash and credit constraints, transfer savings,
and settle payments.  A dedicated seeded random generator drives expenditure
variation, market returns, shocks, and other probabilistic transitions so that
unrelated random draws need not perturb the economy sequence.

Shared sector and government pools make counterparties visible instead of
creating every monetary flow without an account.  Agent-to-agent payments and
some rent or labor flows can be routed to modeled recipients; investment
returns combine a common market component with idiosyncratic variation.
Ledger and conservation checks can reveal implementation defects such as
unbalanced internal transfers.  They cannot establish that the tax schedule,
consumption response, labor process, credit behavior, or macro dynamics is an
empirically valid representation of a real economy.  Accounting consistency
is therefore a software invariant, not a behavioral-validation result.

\subsection{Coupling and ownership boundaries}

Table~\ref{tab:component-coupling} makes the principal shared-state paths
explicit.  ``Implemented'' refers to an inspectable current mechanism;
``partial boundary'' indicates that the mechanism is connected but retains
direct calls, shared dictionaries, or specialized writers rather than a fully
isolated plugin interface.

\begin{table}[t]
\centering
\small
\setlength{\tabcolsep}{3.5pt}
\begin{tabular}{@{}p{0.14\linewidth}p{0.17\linewidth}p{0.17\linewidth}p{0.27\linewidth}p{0.16\linewidth}@{}}
\toprule
Subsystem & Reads & Writes & Agent-visible consequence & Evidence status \\
\midrule
Memory and cognition & Episodes, state, context, retrieved entries & Memories, intentions, habits, relationships & Recalled context, priorities, reflection, and later action bias & Implemented; human realism unvalidated \\
Interests and skills & Profile, episodes, social growth context & Growth profile, private skills & Practice focus and procedural prompt context & Implemented; partial plugin boundary \\
World and local physics & City graph, time, weather, locations & Routes, travel state, occupancy snapshots & Feasible destinations, travel costs, crowding and closure cues & Implemented; direct integration \\
Environment & Configuration, history, optional LLM output & Events, weather and macro context & Interruptions and perceived exogenous conditions & Implemented; generator-mode dependent \\
Social network & Profiles, contact history, interactions & Typed ties, closeness, obligation, tiers & Encounter probability and relationship context & Implemented; mechanism validity unassessed \\
Life events & Persisted event queue, clock & Consumption flag and configured state effects & Scheduled exogenous changes to targeted agents & Implemented; committed plugin-integrated event producer \\
Economy & Activity, balances, calendar, sectors, macro state & Accounts, ledger, budgets, pools, agent state & Income, spending constraints, taxes, credit, returns and payments & Implemented; primarily hook-based \\
Work & Capability cache, activity, growth, market jobs & Queue results, artifacts, settlement and state effects & Produced files, earnings, and experience records & Implemented; artifact quality unvalidated \\
\bottomrule
\end{tabular}
\caption{Coupling among agent and societal components.  Evidence status
describes implementation visibility, not empirical validity.}
\label{tab:component-coupling}
\end{table}

The table also exposes the current modularity limit.  Interests and skills have
plugin entry points, and economy, local physical perception, spatial
preferences, work, and dynamic behavior now have built-in adapters.  Memory,
environment evolution, and social-network operations still exchange mutable
agent dictionaries or runtime objects directly.  This arrangement is usable
and inspectable, yet replacing
one mechanism can change prompts, state schemas, random draws, and specialized
artifacts elsewhere.  Controlled comparisons should therefore freeze the
complete coupled configuration rather than naming only the focal intervention.
